\section{Bone Scan}


\subsubsection*{What Is This Test?}
A bone scan uses a radioactive tracer that accumulates in areas of high bone turnover. It detects fractures, infections, arthritis, and cancer spread to bones earlier than X-rays. Whole-body scan takes 3-4 hours after injection.

\subsubsection*{Alternative Names}
Nuclear medicine bone scan, bone scintigraphy, technetium-99m bone scan, whole-body bone scan


\subsubsection*{Principle}
Bone scan uses technetium-99m-labeled bisphosphonates (MDP or HDP) injected intravenously. The tracer is taken up by bone in proportion to blood flow and osteoblastic activity. A gamma camera detects the emitted 140 keV gamma photons. Increased uptake (hot spots) indicates increased bone turnover from metastases, infection, or fracture.

\subsubsection*{History}
Nuclear medicine bone scanning was developed in the 1960s using technetium-99m labeled diphosphonates. It became a routine clinical tool in the 1970s for detecting metastatic bone disease.

\subsubsection*{Why This Test Is Ordered}
Your doctor orders a bone scan when they suspect bone metastases (cancer that has spread to your bones), when you have a suspected bone infection, when you have a fracture that does not show up on X-ray, or when evaluating metabolic bone disease. It answers the question: are there areas of abnormal bone activity throughout your skeleton?

\subsubsection*{Is This a Primary or Secondary Test}
A secondary test. Usually ordered after other tests suggest a bone problem, or as part of cancer staging.

\subsubsection*{Bone-scan uptake and its limits}
A bone scan can reveal active remodeling, but uptake does not identify its cause by itself. Performance for metastasis, fracture, or infection must be considered separately and correlated with targeted imaging.

\subsubsection*{How This Test Is Performed}
A small amount of radioactive material is injected into your vein. You wait three to four hours while the material distributes to your bones. Then you lie on a table while a special camera takes images of your entire skeleton. The test takes thirty to sixty minutes. There is no pain involved.

\subsubsection*{What Preparation Is Needed}
No special preparation is needed. Tell your doctor if you are pregnant or breastfeeding.

\subsubsection*{Contraindications and Precautions}
Pregnancy and breastfeeding require examination-specific risk assessment, dose optimization, and current radiopharmaceutical guidance. For FDG studies, follow the protocol for preparation and glucose management.

\subsubsection*{Risks}
Risks associated with nuclear medicine include:
\begin{itemize}
    \item \textbf{Radiation exposure}: Low-dose radiation exposure from the administered radiopharmaceutical tracer.
    \item \textbf{Injection site discomfort}: Minor pain, bruising, or rare infection at the intravenous injection site.
\end{itemize}


You can resume normal activities immediately. The radioactive material will naturally decay and leave your body within a day or two. Drink plenty of fluids to help flush it out. Results typically take one to three days.

\subsubsection*{Understanding Your Results}

\subsubsection*{Normal Results}
A normal result shows symmetric tracer uptake throughout the skeleton with no areas of abnormally increased or decreased activity.

\begin{center}
\begin{tabular}{ll}
\toprule
\textbf{Population} & \textbf{Reference Range} \\
\midrule
Adult Female & Symmetric tracer uptake \\
Adult Male & Symmetric tracer uptake \\
Children & Apply age-specific imaging and safety criteria \\
Elderly & Apply age- and comorbidity-specific criteria \\
Pregnancy & Apply pregnancy-specific justification and safety criteria \\
\bottomrule
\end{tabular}
\end{center}
\clearpage
